\chapter{Le système immunitaire}

\noindent\textit{Face à un microbe déjà rencontré, l'organisme réagit plus vite et plus
fort que face à un microbe nouveau. Cette « mémoire » repose sur des cellules et des
molécules spécialisées, que ce chapitre présente, avant d'expliquer le principe de la
vaccination.}
\vspace{8pt}

\connecteBanner

\section{Les organes du système immunitaire}

\begin{propriete}[Organes lymphoïdes]
La \textbf{moelle osseuse} produit toutes les cellules du sang, dont les cellules
immunitaires. Le \textbf{thymus} est le lieu de maturation des lymphocytes T. Les
\textbf{ganglions lymphatiques} et la \textbf{rate} sont des lieux où les lymphocytes
rencontrent les antigènes et se multiplient en cas d'infection (d'où le gonflement des
ganglions lors d'une infection).
\end{propriete}

\section{Antigènes, lymphocytes et anticorps}

\begin{definition}[Antigène]
Un \textbf{antigène} est une molécule étrangère (portée par un microbe, par exemple) que
le système immunitaire est capable de \textbf{reconnaître}, déclenchant une réponse
immunitaire spécifique dirigée contre elle.
\end{definition}

\begin{propriete}[Deux types de lymphocytes]
Les \textbf{lymphocytes B}, une fois activés, se transforment en \textbf{plasmocytes} qui
\textbf{produisent des anticorps} spécifiques de l'antigène rencontré. Les
\textbf{lymphocytes T} agissent directement : certains détruisent les cellules infectées,
d'autres coordonnent la réponse immunitaire.
\end{propriete}

\begin{definition}[Anticorps]
Un \textbf{anticorps} est une molécule produite par les plasmocytes, \textbf{spécifique
d'un antigène précis} : il se fixe sur cet antigène (comme une clé dans une serrure), ce
qui le neutralise ou facilite sa destruction par phagocytose.
\end{definition}

\begin{illus}
\begin{tikzpicture}[scale=0.85]
\filldraw[onbo!25,draw=onbo,line width=1.2pt] (0,0) circle(0.6);
\node[font=\small,align=center] at (0,-1.1) {antigène};
\draw[onbgreen,line width=1.5pt] (2.6,0.5) -- (3.4,0.5) -- (3.4,-0.5) -- (2.6,-0.5);
\node[font=\small,align=center] at (3,-1.1) {anticorps};
\draw[<->,onbmuted,line width=1pt,dashed] (0.7,0) -- (2.45,0);
\node[above,font=\scriptsize,onbmuted] at (1.55,0.35) {liaison spécifique};
\end{tikzpicture}
\fig{Figure — L'anticorps se lie spécifiquement à l'antigène qui a provoqué sa production.}
\end{illus}

\section{Réponse primaire et réponse secondaire}

\begin{propriete}[Première rencontre : réponse primaire]
Lors d'un \textbf{premier contact} avec un antigène, la production d'anticorps est
\textbf{lente} et le taux d'anticorps atteint reste \textbf{modéré}. Une partie des
lymphocytes activés deviennent des \textbf{lymphocytes mémoire}, qui persistent longtemps
dans l'organisme.
\end{propriete}

\begin{propriete}[Deuxième rencontre : réponse secondaire]
Lors d'un \textbf{contact ultérieur} avec le \textbf{même} antigène, grâce aux
lymphocytes mémoire, la réponse est \textbf{beaucoup plus rapide} et le taux d'anticorps
atteint un niveau \textbf{plus élevé} : c'est la \textbf{mémoire immunitaire}.
\end{propriete}

\begin{illus}
\begin{tikzpicture}[scale=0.8]
\draw[->,onbmuted] (0,0) -- (10.5,0) node[right,font=\small] {temps};
\draw[->,onbmuted] (0,0) -- (0,5.2) node[above,font=\small] {taux d'anticorps};
\draw[onbo,line width=1.5pt,smooth]
  plot coordinates {(1,0.1) (1.5,0.3) (2.2,1.1) (3,1.4) (4,1.1) (5,0.6)};
\node[onbo,font=\scriptsize] at (3,1.8) {réponse primaire};
\draw[onbgreen,line width=1.5pt,smooth]
  plot coordinates {(6,0.1) (6.3,0.5) (6.8,3.2) (7.4,4.6) (8.2,4.3) (9.5,2.8)};
\node[onbgreen,font=\scriptsize] at (7.6,5) {réponse secondaire};
\draw[onbmuted,dashed] (1,0) -- (1,0.1);
\node[below,font=\scriptsize] at (1,-0.1) {1\textsuperscript{er} contact};
\draw[onbmuted,dashed] (6,0) -- (6,0.1);
\node[below,font=\scriptsize] at (6,-0.1) {2\textsuperscript{e} contact};
\end{tikzpicture}
\fig{Figure — Réponse primaire (lente, modérée) et réponse secondaire (rapide, intense) au même antigène.}
\end{illus}

\section{Le principe de la vaccination}

\begin{definition}[Vaccination]
La \textbf{vaccination} consiste à introduire dans l'organisme un antigène
\textbf{atténué} ou \textbf{inactivé} (le vaccin), incapable de provoquer la maladie, mais
capable de déclencher une \textbf{réponse primaire} et de créer des \textbf{lymphocytes
mémoire}.
\end{definition}

\begin{propriete}[Intérêt de la vaccination]
En cas de rencontre ultérieure avec le \textbf{véritable} agent pathogène, l'organisme
vacciné déclenche immédiatement une \textbf{réponse secondaire}, rapide et efficace, qui
élimine le microbe avant qu'il ne provoque la maladie.
\end{propriete}

\begin{exemple}
Un enfant vacciné contre la rougeole développe des lymphocytes mémoire spécifiques du
virus de la rougeole ; s'il est un jour exposé au virus réel, sa réponse immunitaire sera
assez rapide pour éviter la maladie.
\end{exemple}

% ═══════════════════════════════════════════════════════════════
\section{L'essentiel à retenir}

\begin{synthese}
\textbf{Antigène.} Molécule étrangère reconnue par le système immunitaire.\\[4pt]
\textbf{Lymphocytes B.} Deviennent des plasmocytes, produisent des anticorps.
\textbf{Lymphocytes T.} Détruisent les cellules infectées, coordonnent la réponse.\\[4pt]
\textbf{Anticorps.} Spécifique d'un antigène précis.\\[4pt]
\textbf{Réponse primaire.} Lente, modérée. \textbf{Réponse secondaire.} Rapide, intense
(grâce aux lymphocytes mémoire).\\[4pt]
\textbf{Vaccination.} Antigène atténué/inactivé $\to$ réponse primaire $+$ mémoire, sans
provoquer la maladie.\\[4pt]
\textbf{Piège fréquent.} Un anticorps est spécifique d'\textbf{un seul} antigène : un
vaccin contre une maladie ne protège pas contre une autre maladie.
\end{synthese}

% ═══════════════════════════════════════════════════════════════
\section{Méthodes \& savoir-faire}

\methode{Distinguer lymphocytes B et lymphocytes T}
Se rappeler : B $\to$ anticorps (via les plasmocytes) ; T $\to$ action directe ou
coordination.
\begin{exemple}
La production d'anticorps contre un virus implique des lymphocytes B.
\end{exemple}

\methode{Lire un graphique de réponse primaire / secondaire}
Comparer la rapidité de montée et le niveau maximal atteint entre le premier et le second
contact avec l'antigène.
\begin{exemple}
Un pic d'anticorps plus élevé et plus rapide lors du second contact illustre la mémoire
immunitaire.
\end{exemple}

\methode{Expliquer le principe d'un vaccin}
Relier l'injection du vaccin (antigène atténué) à la réponse primaire, puis la mémoire
immunitaire à la réponse secondaire rapide en cas d'infection réelle.
\begin{exemple}
Un rappel de vaccin renforce la mémoire immunitaire déjà constituée par les injections
précédentes.
\end{exemple}

% ═══════════════════════════════════════════════════════════════
\section{Exercices d'application}
\begin{multicols}{2}

\rubrique{Organes et cellules}
\exo{}\diff{1} Où sont produites les cellules du système immunitaire ?

\exo{}\diff{2} Quel est le rôle du thymus ?

\exo{}\diff{1} Qu'est-ce qu'un antigène ?

\rubrique{Lymphocytes et anticorps}
\exo{}\diff{2} Que deviennent les lymphocytes B activés ?

\exo{}\diff{2} Qu'est-ce qu'un anticorps ?

\exo{}\diff{2} Quel est le rôle des lymphocytes T ?

\rubrique{Réponse primaire et secondaire}
\exo{}\diff{2} Quelle est la différence entre réponse primaire et réponse secondaire ?

\exo{}\diff{2} Que deviennent certains lymphocytes après un premier contact avec un
antigène ?

\rubrique{Vaccination}
\exo{}\diff{1} Qu'est-ce qu'un vaccin ?

\exo{}\diff{2} Pourquoi un vaccin ne provoque-t-il normalement pas la maladie qu'il
prévient ?

% ═══════════════════════════════════════════════════════════════
\end{multicols}

\section{Exercices d'entraînement}
\begin{multicols}{2}

\rubrique{Organes et cellules}
\exo{}\diff{3} Pourquoi les ganglions lymphatiques gonflent-ils souvent lors d'une
infection ?

\rubrique{Lymphocytes et anticorps}
\exo{}\diff{3} Pourquoi dit-on qu'un anticorps est « spécifique » d'un antigène ?

\exo{}\diff{2} Un antigène nouveau apparaît dans l'organisme. Quelle cellule va produire
les anticorps correspondants, après activation ?

\rubrique{Réponse primaire et secondaire}
\exo{}\diff{3} À l'aide d'un graphique donnant le taux d'anticorps en fonction du temps
lors de deux contacts avec le même antigène, comment reconnaître la courbe de la réponse
secondaire ?

\exo{}\diff{3} Pourquoi la réponse secondaire est-elle plus rapide que la réponse
primaire ?

\rubrique{Vaccination}
\exo{}\diff{3} Expliquer pourquoi un rappel de vaccin est parfois nécessaire plusieurs
années après la première injection.

\exo{}\diff{3} Un antigène atténué contenu dans un vaccin est-il capable de déclencher
une réponse immunitaire ? Justifier.

\exo{}\diff{2} Pourquoi ne peut-on pas se protéger de toutes les maladies avec un seul
vaccin ?

\exo{}\diff{3} Expliquer le lien entre vaccination et réponse secondaire.

% ═══════════════════════════════════════════════════════════════
\end{multicols}

\section{Sujets type BEPC}

\sujetbac[Lymphocytes et anticorps]
\begin{enumerate}[label=\arabic*)]
\item Qu'est-ce qu'un antigène ?
\item Que deviennent les lymphocytes B activés ?
\item Qu'est-ce qu'un anticorps, et comment agit-il ?
\item Quel est le rôle des lymphocytes T ?
\end{enumerate}

\sujetbac[Réponse primaire et secondaire]
\begin{enumerate}[label=\arabic*)]
\item Décris les caractéristiques de la réponse primaire lors d'un premier contact avec
un antigène.
\item Décris les caractéristiques de la réponse secondaire lors d'un second contact avec
le même antigène.
\item Quelles cellules expliquent la différence entre ces deux réponses ?
\item Représente, en une phrase, la différence entre les deux courbes sur un graphique
taux d'anticorps en fonction du temps.
\end{enumerate}

\sujetbac[La vaccination]
\begin{enumerate}[label=\arabic*)]
\item Qu'est-ce qu'un vaccin ?
\item Quelle réponse immunitaire un vaccin déclenche-t-il lors de l'injection ?
\item Que se passe-t-il si la personne vaccinée rencontre plus tard le véritable agent
pathogène ?
\item Pourquoi la vaccination est-elle une mesure de santé publique efficace ?
\end{enumerate}

% ═══════════════════════════════════════════════════════════════
\section{Corrigés}
\begin{multicols}{2}

\corrige{de l'exercice 1}
Dans la moelle osseuse.

\corrige{de l'exercice 2}
La maturation des lymphocytes T.

\corrige{de l'exercice 3}
Une molécule étrangère reconnue par le système immunitaire, qui déclenche une réponse
spécifique.

\corrige{de l'exercice 4}
Ils se transforment en plasmocytes, qui produisent des anticorps.

\corrige{de l'exercice 5}
Une molécule spécifique d'un antigène, qui s'y fixe pour le neutraliser ou faciliter sa
destruction.

\corrige{de l'exercice 6}
Détruire directement les cellules infectées ou coordonner la réponse immunitaire.

\corrige{de l'exercice 7}
La réponse primaire est lente et modérée ; la réponse secondaire est rapide et intense.

\corrige{de l'exercice 8}
Ils deviennent des lymphocytes mémoire.

\corrige{de l'exercice 9}
Une préparation contenant un antigène atténué ou inactivé, destinée à déclencher une
réponse immunitaire protectrice.

\corrige{de l'exercice 10}
Car l'antigène qu'il contient est atténué ou inactivé : il ne peut pas provoquer la
maladie, mais reste capable de déclencher une réponse immunitaire.

\corrige{de l'exercice 11}
Car ils s'y multiplient en réponse à une infection, ce qui augmente leur taille.

\corrige{de l'exercice 12}
Car il se fixe uniquement sur l'antigène qui a provoqué sa production, comme une clé dans
une serrure.

\corrige{de l'exercice 13}
Les lymphocytes B, transformés en plasmocytes.

\corrige{de l'exercice 14}
Elle monte plus rapidement et atteint un niveau plus élevé que la réponse primaire.

\corrige{de l'exercice 15}
Car des lymphocytes mémoire, déjà présents depuis le premier contact, réagissent
immédiatement, sans le délai nécessaire à leur formation initiale.

\corrige{du sujet 1}
1) Une molécule étrangère reconnue par le système immunitaire.\\
2) Ils se transforment en plasmocytes, producteurs d'anticorps.\\
3) Une molécule spécifique d'un antigène qui s'y fixe pour le neutraliser ou faciliter
sa destruction.\\
4) Détruire directement les cellules infectées ou coordonner la réponse immunitaire.

\corrige{du sujet 2}
1) Elle est lente et le taux d'anticorps atteint reste modéré.\\
2) Elle est rapide et le taux d'anticorps atteint un niveau plus élevé.\\
3) Les lymphocytes mémoire, formés lors du premier contact.\\
4) La courbe de la réponse secondaire monte plus vite et plus haut que celle de la
réponse primaire.

\corrige{du sujet 3}
1) Une préparation contenant un antigène atténué ou inactivé, capable de déclencher une
réponse immunitaire protectrice sans provoquer la maladie.\\
2) Une réponse primaire, avec formation de lymphocytes mémoire.\\
3) Une réponse secondaire, rapide et efficace, se déclenche et élimine le microbe avant
qu'il ne provoque la maladie.\\
4) Car elle permet de protéger un grand nombre de personnes contre des maladies graves,
sans leur faire courir le risque de la maladie elle-même.

\end{multicols}
\exercicesBanner
